\paragraph{Using an already compiled main.}
When the native main-only executable is already available to the instrumentation client, we remove exactly that component's measured build charge from each scoped choice. We still charge the observer build, any additional dispatch, and the repeated call time. At $N=1$, the scoped rule then costs 0.689 times native joint across the 24 cells; at $N=100$, 0.862; at $N=1000$, 0.995; and at $N=10000$, 1.053. The corresponding cold-build ratios are 1.348, 1.285, 1.171 and 1.083. Reusing a main therefore changes the short-lived deployment decision. For a cold build on this panel, native joint has the lowest aggregate cost at all four budgets. For a client with a reusable executable serving very few calls, the fixed scoped alternative merits comparison. Possessing only an HLO listing yields no executable-build credit.

A single amortization threshold would obscure differences among graphs. For the small independent-input graph on one allowed CPU, the scoped choice first beats cold native joint after 116--236 calls across the five processes. At default affinity none of its five paired build-and-call measurements has a positive integer break-even point. The large row-centered graph also has no cold break-even in any of the five default-affinity processes. The artifact gives the full intervals, including finite winning intervals when a cheaper build buys a slower executable.
